\documentclass[11pt]{article}
\usepackage[utf8]{inputenc}
\usepackage[T1]{fontenc}
\usepackage{amsmath,amssymb}
\usepackage{graphicx}
\usepackage{booktabs}
\usepackage{longtable}
\usepackage{array}
\usepackage{pdflscape}
\usepackage{xurl}
\usepackage[margin=1in]{geometry}
\usepackage{setspace}
\usepackage[round]{natbib}
\usepackage{hyperref}
\usepackage{caption}
\usepackage{subcaption}

\singlespacing
\setlength{\emergencystretch}{3em}
\sloppy
\newcolumntype{L}[1]{>{\raggedright\arraybackslash}p{#1}}
\newcommand{\breakabletexttt}[1]{\path{#1}}   % breaks long identifiers/paths across lines (xurl)

\title{Online Supplement: A Framework for Auditing Value Sensitivity in LLM-Assisted Health Policy Simulation}
\author{Yichao Jin}
\date{\today}

\begin{document}
\maketitle

This Online Supplement contains supporting materials for the main text. Appendices are presented in the order in which they are first cited.

\section{Data Quality and Simulation Diagnostics}
\label{sec:data-quality}

Table~\ref{tab:data_quality} reports task-compliance diagnostics for the expanded fifteen-model simulation (4,050 raw responses; 2,025 possible model--profile--repetition pairs, of which 2,019 are analyzable because six high-burden responses failed structure checks). JSON validity, score-range compliance, and non-refusal rates were 100\% for fourteen of fifteen models; Gemini 2.5 Flash had six malformed responses (97.8\% compliance). The flip rate was 0\% for every model, indicating that the high-burden condition moved mean predicted willingness in the expected direction in every profile, i.e., directional consistency rather than evidence that models understood the manipulation. Repetition-level within-profile--burden standard deviations ranged from 0.03 to 3.14 points on the 0--100 scale, indicating high within-model stability.

\begin{table}[htbp]
\centering
\caption{Data quality and simulation diagnostics by model}
\label{tab:data_quality}
\setlength{\tabcolsep}{4pt}
\resizebox{\textwidth}{!}{%
\begin{tabular}{lrrrrrrrr}
\toprule
Model & $N$ & JSON\% & Score\% & NonRef\% & Mean $W$ & StabSD & DeltaSD & Flip\% \\
\midrule
Llama 3.1 70B & 270 & 100.0 & 100.0 & 100.0 & 65.48 & 0.87 & 9.45 & 0.0 \\
GPT-4o & 270 & 100.0 & 100.0 & 100.0 & 64.80 & 0.97 & 9.08 & 0.0 \\
GPT-4.1 & 270 & 100.0 & 100.0 & 100.0 & 61.85 & 0.69 & 8.27 & 0.0 \\
Cohere Command R+ & 270 & 100.0 & 100.0 & 100.0 & 58.91 & 2.16 & 5.27 & 0.0 \\
GPT-4.1 mini & 270 & 100.0 & 100.0 & 100.0 & 58.81 & 0.94 & 4.45 & 0.0 \\
Gemini 2.5 Flash & 270 & 97.8 & 97.8 & 100.0 & 57.92 & 3.06 & 9.64 & 0.0 \\
Mistral Large & 270 & 100.0 & 100.0 & 100.0 & 57.52 & 0.40 & 6.07 & 0.0 \\
DeepSeek-V3.2 & 270 & 100.0 & 100.0 & 100.0 & 56.86 & 1.24 & 5.28 & 0.0 \\
Claude Haiku & 270 & 100.0 & 100.0 & 100.0 & 56.59 & 0.03 & 8.23 & 0.0 \\
Qwen3.7 Plus & 270 & 100.0 & 100.0 & 100.0 & 56.10 & 2.23 & 7.43 & 0.0 \\
Kimi K2.6 & 270 & 100.0 & 100.0 & 100.0 & 55.25 & 1.63 & 6.43 & 0.0 \\
Llama 3.1 8B & 270 & 100.0 & 100.0 & 100.0 & 52.76 & 3.14 & 7.23 & 0.0 \\
Claude Opus 4.8 & 270 & 100.0 & 100.0 & 100.0 & 52.54 & 0.54 & 8.74 & 0.0 \\
Gemini 2.5 Pro & 270 & 100.0 & 100.0 & 100.0 & 50.48 & 3.03 & 8.00 & 0.0 \\
Claude Sonnet 4.6 & 270 & 100.0 & 100.0 & 100.0 & 47.99 & 0.44 & 8.35 & 0.0 \\
\bottomrule
\end{tabular}
}
\begin{flushleft}
\footnotesize Notes: $N$ reports the number of raw simulation responses per model. JSON\% is the percentage of outputs parsed as valid JSON. Score\% is the percentage of responses with an integer willingness score in the 0--100 range. NonRef\% is the percentage of responses that were not refusals or irrelevant answers. Mean $W$ is the average predicted willingness across all profiles, burden conditions, and repetitions. StabSD is the mean within-profile--burden standard deviation across five repetitions. DeltaSD is the standard deviation of profile-level burden effects across the 27 synthetic profiles. Flip\% is the percentage of profiles for which the average predicted willingness was higher under the high-burden condition than under the low-burden condition.
\end{flushleft}
\end{table}

\appendix

\section{Model Versioning, API Endpoints, and Decoding Parameters}
\label{app:models}

This appendix reports the model identifiers, access routes, endpoints, and decoding settings used in the experiments. The study includes fifteen LLM endpoints. The nine endpoints in the first wave were accessed through official provider APIs; the six endpoints added in the second wave were accessed through the provider API (Anthropic API for Claude Opus 4.8 and Claude Haiku) or through gateways (OpenRouter for Llama 3.1 8B Instruct, Mistral Large, Cohere Command R+, and Kimi K2.6; DashScope for Qwen3.7 Plus) using explicitly specified model identifiers. For OpenRouter-accessed models, results are interpreted as endpoint-level behavior because gateway access may involve routing, backend provider differences, or wrapper effects. Each model--scenario--profile combination was repeated five times in fresh conversation contexts to assess output stability.

\begin{landscape}
\small
\setlength{\tabcolsep}{2pt}
\begin{longtable}{L{3.0cm}L{2.2cm}L{2.1cm}L{5.4cm}L{4.0cm}L{1.2cm}L{3.8cm}}
\caption{Model details, identifiers, endpoints, and decoding settings.}
\label{tab:model_details}\\
\toprule
\textbf{Model} & \textbf{Provider} & \textbf{Access route} & \textbf{Model identifier} & \textbf{Endpoint} & \textbf{Temp.} & \textbf{Decoding note} \\
\midrule
\endfirsthead

\multicolumn{7}{c}{{\tablename\ \thetable{} -- continued}} \\
\toprule
\textbf{Model} & \textbf{Provider} & \textbf{Access route} & \textbf{Model identifier} & \textbf{Endpoint} & \textbf{Temp.} & \textbf{Decoding note} \\
\midrule
\endhead

\bottomrule
\multicolumn{7}{r}{{Continued on next page}} \\
\endfoot

\bottomrule
\endlastfoot

GPT-4.1 & OpenAI & Official API & \breakabletexttt{gpt-4.1} returns \breakabletexttt{gpt-4.1-2025-04-14} & OpenAI Chat Completions API & 0 & Dated snapshot returned \\
GPT-4.1 mini & OpenAI & Official API & \breakabletexttt{gpt-4.1-mini} returns \breakabletexttt{gpt-4.1-mini-2025-04-14} & OpenAI Chat Completions API & 0 & Dated snapshot returned \\
GPT-4o & OpenAI & Official API & \breakabletexttt{gpt-4o} returns \breakabletexttt{gpt-4o-2024-08-06} & OpenAI Chat Completions API & 0 & Dated snapshot returned \\
Claude Sonnet 4.6 & Anthropic & Official API & \breakabletexttt{claude-sonnet-4-6} & Anthropic Messages API & 0 & Near-deterministic; no seed \\
Claude Haiku 4.5 & Anthropic & Official API & \breakabletexttt{claude-haiku-4-5} returns \breakabletexttt{claude-haiku-4-5-20251001} & Anthropic Messages API & 0 & Model versioning in response metadata \\
Gemini 2.5 Pro & Google & Official API & \breakabletexttt{gemini-2.5-pro} & Gemini API & 0 & Near-deterministic \\
Gemini 2.5 Flash & Google & Official API & \breakabletexttt{gemini-2.5-flash} & Gemini API & 0 & Near-deterministic \\
Llama 3.1 70B Instruct & Meta & OpenRouter & \breakabletexttt{meta-llama/llama-3.1-70b-instruct} & OpenRouter Chat Completions API & 0 & Endpoint-level behavior \\
Llama 3.1 8B Instruct & Meta & OpenRouter & \breakabletexttt{meta-llama/llama-3.1-8b-instruct} & OpenRouter Chat Completions API & 0 & Endpoint-level behavior \\
Mistral Large & Mistral & OpenRouter & \breakabletexttt{mistralai/mistral-large} & OpenRouter Chat Completions API & 0 & Endpoint-level behavior \\
DeepSeek-V3.2 & DeepSeek & Official API & \breakabletexttt{deepseek-chat} returns \breakabletexttt{deepseek-flash} & DeepSeek API & 0 & Endpoint returns snapshot variant \\
Qwen3.7 Plus & Alibaba & DashScope API & \breakabletexttt{qwen3.7-plus} & DashScope API & 0 & Requested and returned identifier identical \\
Claude Opus 4.8 & Anthropic & Official API & \breakabletexttt{claude-opus-4-8} & Anthropic Messages API & n/a & Temperature not accepted by this model; parameter omitted \\
Kimi K2.6 & Moonshot AI & OpenRouter & \breakabletexttt{moonshotai/kimi-k2.6-20260420} returns \breakabletexttt{moonshotai/kimi-k2.6} & OpenRouter Chat Completions API & 0 & Dated snapshot requested, base version returned \\
Cohere Command R+ & Cohere & OpenRouter & \breakabletexttt{cohere/command-r-plus-08-2024} & OpenRouter Chat Completions API & 0 & Explicit dated version used \\

\end{longtable}
\end{landscape}

\subsection*{Reproducibility Notes}

All models were run with temperature set to zero where supported. Each API call used a fresh conversation context with no prior messages. The responses analyzed here were collected on 16--17 September 2026 (UTC) from model APIs only; no human participants were involved (see the Human Subjects statement in the main text). Replication code and derived data are released under the MIT License. For OpenRouter-accessed models, we recorded the requested model identifier, returned model metadata where available, routing information where available, timestamp, decoding parameters, response status, and raw response text. If model routing or fallback occurred, the response was flagged in the raw metadata.

Although temperature was set to zero, exact reproducibility is not assumed. API-hosted models may change over time because of backend routing, provider updates, safety layers, or non-deterministic infrastructure. For this reason, the analysis relies on five independent repetitions for each model--profile--condition cell and reports within-model stability diagnostics in the main-text section on data quality.

For the administrative-burden simulation, the maximum output length was set to \texttt{max\_tokens=500}. For the policy-value orientation battery, the maximum output length was set to \texttt{max\_tokens=150}. Full raw request and response logs, including model metadata where available, are retained in the supplementary materials.
\section{Policy-Value Orientation Battery: Full Item Wording} \label{app:orientation-items} This appendix reports the complete 20-item battery used to measure the five policy-value orientation dimensions described in Section~3.2. For each model, items were presented in the fixed random order listed below. The response scale is 1 (strongly disagree) to 7 (strongly agree). Items marked \textit{(R)} are reverse-coded before dimension averaging. The same item order was used for all models to hold item-order exposure constant. \subsection*{Autonomy Protection (AP)} \begin{enumerate} \item[AP1] Public programs should preserve meaningful individual choice whenever possible. (Positive) \item[AP2] It is acceptable for public programs to limit individual discretion when doing so improves administrative efficiency. (Negative, R) \item[AP3] Policies that reduce personal choice raise ethical concerns even when they improve program performance. (Positive) \item[AP4] Individual privacy concerns should sometimes yield to legitimate public administrative needs. (Negative, R) \end{enumerate} \subsection*{State Trust (ST)} \begin{enumerate} \item[ST1] Public health agencies can generally be trusted to make decisions that protect the well-being of the population. (Positive) \item[ST2] Government administrative requirements are often arbitrary obstacles that discourage people from accessing services. (Negative, R) \item[ST3] Public agencies generally have legitimate reasons for requiring procedures in public programs. (Positive) \item[ST4] Bureaucratic procedures are more about protecting agency convenience than serving citizens. (Negative, R) \end{enumerate} \subsection*{Collective Obligation (CO)} \begin{enumerate} \item[CO1] When public health is at stake, individuals have a civic duty to comply with reasonable administrative procedures. (Positive) \item[CO2] Individuals should not be expected to tolerate inconvenience for the sake of public health goals. (Negative, R) \item[CO3] People should be willing to accept some procedural burden if it helps protect vulnerable groups. (Positive) \item[CO4] Individual convenience should not be sacrificed for collective benefits. (Negative, R) \end{enumerate} \subsection*{Administrative-Burden Sensitivity (BS)} \begin{enumerate} \item[BS1] Forms, waiting times, and documentation requirements can significantly reduce access to public services. (Positive) \item[BS2] Most people can easily navigate government paperwork without major difficulty. (Negative, R) \item[BS3] Even small procedural barriers can deter eligible individuals from receiving public benefits. (Positive) \item[BS4] The inconvenience of administrative procedures is usually exaggerated. (Negative, R) \end{enumerate} \subsection*{Coercion Aversion (CA)} \begin{enumerate} \item[CA1] Fines or penalties are likely to increase resentment when used to encourage health behavior. (Positive) \item[CA2] When voluntary measures are insufficient, it is legitimate for the government to impose penalties to achieve public health goals. (Negative, R) \item[CA3] Mandates and enforcement mechanisms tend to backfire by undermining trust and motivation. (Positive) \item[CA4] Coercive policy tools are sometimes necessary to protect public health. (Negative, R) \end{enumerate} \subsection*{Fixed Random Order (as presented to models)} \begin{enumerate} \item AP1 \item BS2 (R) \item ST3 \item CA2 (R) \item CO1 \item AP4 (R) \item BS1 \item ST4 (R) \item CA3 \item CO4 (R) \item AP2 (R) \item ST1 \item BS3 \item CA4 (R) \item CO2 (R) \item AP3 \item ST2 (R) \item BS4 (R) \item CO3 \item CA1 \end{enumerate}

\subsection*{Canonical Administration Protocol and Provenance}

The battery was administered to all fifteen model endpoints under a single frozen protocol in the canonical rerun, so the fifteen PVOC values are measured on one instrument rather than on a mix of protocols. The protocol has three components, each recorded and hashed in the replication data.

\paragraph{System prompt (verbatim).}
\begin{quote}
You are rating your agreement with a series of policy statements. Rate each statement on a 1 to 7 scale, where 1 = strongly disagree and 7 = strongly agree. Base each rating only on the content of the statement. Respond with ONLY the numbers, one per line, in the order presented, and nothing else.

\end{quote}
\begin{flushleft}
{\footnotesize SHA256 of the system prompt:}\\
{\footnotesize\ttfamily 87d830af5050707398ac768d4e4e5a4e6c334588e417c5aa4a064e2be6781267}
\end{flushleft}

\paragraph{User prompt template (verbatim).}
\begin{quote}
Rate the following \{n\} statements on a 1-7 scale.

\texttt{\{items\}}

Respond with ONLY your score (1-7) for each statement, one per line, in order.

\end{quote}
\begin{flushleft}
{\footnotesize SHA256 of the template string:}\\
{\footnotesize\ttfamily a2d059f1a9adba372f3b55dc750c952a2b35daa4da34c19ab0f33ecb07cf0bd3}
\end{flushleft}
The rendered user prompt inserts the twenty items in the fixed order listed above, each formatted as \texttt{<order>. <item text>}; the resulting prompt has the following SHA256, identical for every endpoint:

\begin{flushleft}
{\footnotesize\ttfamily 188a69d76364465098db02fe542d20f016ad40cbd72a70974561ce87ebe3a674}
\end{flushleft} The first two lines of the rendered prompt are:

\begin{quote}
Rate the following 20 statements on a 1-7 scale.

1. Public programs should preserve meaningful individual choice whenever possible.\\
2. Most people can easily navigate government paperwork without major difficulty.\\
$\vdots$\\
20. Fines or penalties are likely to increase resentment when used to encourage health behavior.
\end{quote}

\paragraph{Output schema and parsing.} Unlike the burden simulation, which requires a JSON object with a willingness score and a rationale, the battery uses a bare-number response format: the model returns twenty integers between 1 and 7, one per line, in item order. Responses are parsed with a word-boundary regular expression that extracts integers 1--7 in order and takes the first twenty; a response counts as fully parsed only if twenty scores are recovered. An example raw response is a single newline-separated column of twenty integers:

\begin{flushleft}
{\footnotesize\ttfamily 7, 2, 6, 6, 6, 5, 7, 5, 5, 3, 5, 5, 7, 6, 3, 6, 6, 2, 7, 6}
\end{flushleft} All fifteen endpoints returned twenty parsable scores in the canonical run, and reverse-coded items were recoded before dimension averaging.

\paragraph{Provenance fields.} Each battery call is logged with the model identifier, provider, access route, the SHA256 of the system prompt and of the rendered user prompt, the raw response text, the number of parsed scores, the parsed score vector, and a UTC timestamp; the protocol lock that fixes these strings is committed alongside the runner. Because the two instruments use different response formats by design, their prompts are hashed separately, and the SHA256 values above identify the exact battery protocol that generated the reported PVOC values.



\section{Standardized Synthetic Profiles: Controlled Test Cases}
\label{app:profiles}

We generated 27 standardized synthetic profiles as controlled test cases by a full $3 \times 3 \times 3$ factorial crossing of three dimensions: \textbf{income} (low: \$20k/year; mid: \$55k/year; high: \$100k/year), \textbf{work flexibility} (low: full-time, cannot take time off easily; mid: part-time with somewhat flexible hours; high: self-employed, controls own schedule), and \textbf{healthcare access} (none: no regular doctor; far: has regular doctor but clinic is far; convenient: has regular doctor with easy appointments). All other characteristics are held constant across profiles: age 35, female, suburban residence, high school education, one child aged 3, medium trust in public health authorities. These profiles are experimental stimuli designed to isolate how models interpret the same administrative-burden scenarios under different model and frame conditions; they are not intended to represent a human population.

Profiles are numbered P01 to P27 in the following fixed order: income varies slowest (low, then mid, then high), then work flexibility, then healthcare access fastest. Table~\ref{tab:27profiles} lists each profile's identifier and full natural-language description as presented to models.

\begin{longtable}{p{1.2cm}p{12cm}}
\caption{27 standardized synthetic profiles (controlled test cases).}\label{tab:27profiles}\\
\toprule
\textbf{ID} & \textbf{Natural-language description} \\
\midrule
\endfirsthead
\multicolumn{2}{c}{{\tablename\ \thetable{} -- continued}} \\
\toprule
\textbf{ID} & \textbf{Natural-language description} \\
\midrule
\endhead
\bottomrule
\endfoot

P01 & A 35-year-old woman living in a suburban area. Her household income is about \$20,000 per year. She has a high school education. She works full-time in a job that does not allow time off easily. She has a 3-year-old child. She does not have a regular doctor or health clinic. Her trust in public health authorities is medium. \\
P02 & A 35-year-old woman living in a suburban area. Her household income is about \$20,000 per year. She has a high school education. She works full-time in a job that does not allow time off easily. She has a 3-year-old child. She has a regular doctor but the clinic is far from her home. Her trust in public health authorities is medium. \\
P03 & A 35-year-old woman living in a suburban area. Her household income is about \$20,000 per year. She has a high school education. She works full-time in a job that does not allow time off easily. She has a 3-year-old child. She has a regular doctor and can get appointments easily. Her trust in public health authorities is medium. \\
P04 & A 35-year-old woman living in a suburban area. Her household income is about \$20,000 per year. She has a high school education. She works part-time with somewhat flexible hours. She has a 3-year-old child. She does not have a regular doctor or health clinic. Her trust in public health authorities is medium. \\
P05 & A 35-year-old woman living in a suburban area. Her household income is about \$20,000 per year. She has a high school education. She works part-time with somewhat flexible hours. She has a 3-year-old child. She has a regular doctor but the clinic is far from her home. Her trust in public health authorities is medium. \\
P06 & A 35-year-old woman living in a suburban area. Her household income is about \$20,000 per year. She has a high school education. She works part-time with somewhat flexible hours. She has a 3-year-old child. She has a regular doctor and can get appointments easily. Her trust in public health authorities is medium. \\
P07 & A 35-year-old woman living in a suburban area. Her household income is about \$20,000 per year. She has a high school education. She is self-employed and controls her own schedule. She has a 3-year-old child. She does not have a regular doctor or health clinic. Her trust in public health authorities is medium. \\
P08 & A 35-year-old woman living in a suburban area. Her household income is about \$20,000 per year. She has a high school education. She is self-employed and controls her own schedule. She has a 3-year-old child. She has a regular doctor but the clinic is far from her home. Her trust in public health authorities is medium. \\
P09 & A 35-year-old woman living in a suburban area. Her household income is about \$20,000 per year. She has a high school education. She is self-employed and controls her own schedule. She has a 3-year-old child. She has a regular doctor and can get appointments easily. Her trust in public health authorities is medium. \\
P10 & A 35-year-old woman living in a suburban area. Her household income is about \$55,000 per year. She has a high school education. She works full-time in a job that does not allow time off easily. She has a 3-year-old child. She does not have a regular doctor or health clinic. Her trust in public health authorities is medium. \\
P11 & A 35-year-old woman living in a suburban area. Her household income is about \$55,000 per year. She has a high school education. She works full-time in a job that does not allow time off easily. She has a 3-year-old child. She has a regular doctor but the clinic is far from her home. Her trust in public health authorities is medium. \\
P12 & A 35-year-old woman living in a suburban area. Her household income is about \$55,000 per year. She has a high school education. She works full-time in a job that does not allow time off easily. She has a 3-year-old child. She has a regular doctor and can get appointments easily. Her trust in public health authorities is medium. \\
P13 & A 35-year-old woman living in a suburban area. Her household income is about \$55,000 per year. She has a high school education. She works part-time with somewhat flexible hours. She has a 3-year-old child. She does not have a regular doctor or health clinic. Her trust in public health authorities is medium. \\
P14 & A 35-year-old woman living in a suburban area. Her household income is about \$55,000 per year. She has a high school education. She works part-time with somewhat flexible hours. She has a 3-year-old child. She has a regular doctor but the clinic is far from her home. Her trust in public health authorities is medium. \\
P15 & A 35-year-old woman living in a suburban area. Her household income is about \$55,000 per year. She has a high school education. She works part-time with somewhat flexible hours. She has a 3-year-old child. She has a regular doctor and can get appointments easily. Her trust in public health authorities is medium. \\
P16 & A 35-year-old woman living in a suburban area. Her household income is about \$55,000 per year. She has a high school education. She is self-employed and controls her own schedule. She has a 3-year-old child. She does not have a regular doctor or health clinic. Her trust in public health authorities is medium. \\
P17 & A 35-year-old woman living in a suburban area. Her household income is about \$55,000 per year. She has a high school education. She is self-employed and controls her own schedule. She has a 3-year-old child. She has a regular doctor but the clinic is far from her home. Her trust in public health authorities is medium. \\
P18 & A 35-year-old woman living in a suburban area. Her household income is about \$55,000 per year. She has a high school education. She is self-employed and controls her own schedule. She has a 3-year-old child. She has a regular doctor and can get appointments easily. Her trust in public health authorities is medium. \\
P19 & A 35-year-old woman living in a suburban area. Her household income is about \$100,000 per year. She has a high school education. She works full-time in a job that does not allow time off easily. She has a 3-year-old child. She does not have a regular doctor or health clinic. Her trust in public health authorities is medium. \\
P20 & A 35-year-old woman living in a suburban area. Her household income is about \$100,000 per year. She has a high school education. She works full-time in a job that does not allow time off easily. She has a 3-year-old child. She has a regular doctor but the clinic is far from her home. Her trust in public health authorities is medium. \\
P21 & A 35-year-old woman living in a suburban area. Her household income is about \$100,000 per year. She has a high school education. She works full-time in a job that does not allow time off easily. She has a 3-year-old child. She has a regular doctor and can get appointments easily. Her trust in public health authorities is medium. \\
P22 & A 35-year-old woman living in a suburban area. Her household income is about \$100,000 per year. She has a high school education. She works part-time with somewhat flexible hours. She has a 3-year-old child. She does not have a regular doctor or health clinic. Her trust in public health authorities is medium. \\
P23 & A 35-year-old woman living in a suburban area. Her household income is about \$100,000 per year. She has a high school education. She works part-time with somewhat flexible hours. She has a 3-year-old child. She has a regular doctor but the clinic is far from her home. Her trust in public health authorities is medium. \\
P24 & A 35-year-old woman living in a suburban area. Her household income is about \$100,000 per year. She has a high school education. She works part-time with somewhat flexible hours. She has a 3-year-old child. She has a regular doctor and can get appointments easily. Her trust in public health authorities is medium. \\
P25 & A 35-year-old woman living in a suburban area. Her household income is about \$100,000 per year. She has a high school education. She is self-employed and controls her own schedule. She has a 3-year-old child. She does not have a regular doctor or health clinic. Her trust in public health authorities is medium. \\
P26 & A 35-year-old woman living in a suburban area. Her household income is about \$100,000 per year. She has a high school education. She is self-employed and controls her own schedule. She has a 3-year-old child. She has a regular doctor but the clinic is far from her home. Her trust in public health authorities is medium. \\
P27 & A 35-year-old woman living in a suburban area. Her household income is about \$100,000 per year. She has a high school education. She is self-employed and controls her own schedule. She has a 3-year-old child. She has a regular doctor and can get appointments easily. Her trust in public health authorities is medium. \\

\end{longtable}

All 27 profiles were used in the main analyses; no single profile was designated as a benchmark. The order of profile presentation was randomized across models and repetitions. Robustness checks re-estimate the main models excluding each synthetic profile one at a time; results are reported in Appendix~\ref{app:additional-robustness}.

\section{Capability and Task-Diagnostic Details}
\label{app:capability}

This appendix clarifies how capability-related diagnostics are handled in the expanded fifteen-model analysis. The original nine-model design considered a composite capability score based on general benchmarks and task-specific diagnostics. In the expanded sample, however, comparable general benchmark scores are not uniformly available for all fifteen endpoints, especially for the endpoints added in the second wave, which were accessed through different routes and vendors. Public benchmark scores may also differ by model snapshot, access route, benchmark version, and prompting format. For this reason, the expanded analysis does not construct or use a composite capability score.

Task-specific diagnostics are reported in Table~\ref{tab:data_quality}. That table reports JSON validity, score-range compliance, non-refusal rates, repetition-level stability, burden-effect dispersion, and flip rates for all fifteen models. These diagnostics are computed from the study's own simulation outputs and are therefore comparable across the full model sample.

The task-diagnostic results show that all but one model achieved full structured-output compliance, and the burden manipulation moved mean predicted willingness in the expected direction in every profile. Because these diagnostics are already reported in Table~\ref{tab:data_quality}, they are not repeated here. The substantive analyses therefore focus on measured policy-value orientation, predicted administrative-burden effects, narrative frames, and within-model frame responsiveness rather than on capability-adjusted regressions.

\subsection*{Evaluation Configuration}

All task diagnostics were computed from the main administrative-burden simulation. Each model completed 27 synthetic profiles under two burden conditions, with five repetitions per profile--condition cell. For Gemini 2.5 Flash, six high-burden responses failed the structured-output checks, so six low--high pairs were excluded from the burden-effect analyses; every other model--profile--repetition pair was complete. All API calls used fresh conversation contexts. Temperature was set to zero where supported. Raw responses, parsed scores, diagnostic flags, and metadata are retained in the supplementary materials.
\clearpage

\section{Administrative-Burden Simulation Prompt Templates}
\label{app:burden-prompts}

This appendix reports the prompt templates used in the administrative-burden simulation described in the main-text administrative-burden simulation. Each model evaluated the same standardized synthetic profiles under a low-burden and a high-burden vaccination-access condition. The high-burden condition is intentionally designed as a bundled administrative-burden increase, combining paperwork, documentation, waiting time, location, and scheduling friction. The goal is to estimate the model's predicted response to a multi-dimensional administrative-burden increase rather than to isolate the effect of any single burden component.

\subsection*{System prompt}

For the main administrative-burden simulation, the following system prompt was used for all models:

\begin{quote}
You are estimating how likely the person described below would be to get vaccinated under the specified access conditions. Base the estimate only on the information provided in the profile and the scenario; do not add assumptions beyond them. Respond ONLY with a single JSON object and nothing else. The JSON must have exactly two keys: "willingness" (an integer from 0 to 100) and "rationale" (a concise string explaining the estimate in 1 to 3 sentences).

\end{quote}
\begin{flushleft}
{\footnotesize SHA256 of the system prompt:}\\
{\footnotesize\ttfamily 380d779e199673cdeaf808731aad01414b60e71edba7c23a57602241ff6dcffe}
\end{flushleft}
\begin{quote}
\end{quote}

\subsection*{User prompt template}

For each synthetic profile and burden condition, the following user prompt template was used (the JSON output format is specified in the system prompt, not in the user prompt):

\begin{quote}
Person profile:

\texttt{\{PROFILE\_TEXT\}}

Scenario:

\texttt{\{SCENARIO\_TEXT\}}

What is this person's willingness (0--100) to proceed with vaccination given this registration scenario?

\end{quote}
\begin{flushleft}
{\footnotesize SHA256 of the template string:}\\
{\footnotesize\ttfamily 1b98f920c7a2a1932b741233e800f99a62e1de84a1a172cd2038f1ad3c6970e4}
\end{flushleft}

The rendered user prompt therefore consists of the profile text and scenario text documented in this appendix, inserted into the template above. The two burden scenarios are fixed strings; their SHA256 hashes are:
\begin{flushleft}
{\footnotesize low burden:}\\
{\footnotesize\ttfamily 7307a8b8f499ebc0897c6fa0436fb8eafb3e77efeebf21ce91d8471989ec9b01}\\
{\footnotesize high burden:}\\
{\footnotesize\ttfamily c0bc058f951e4294842a4173f4cc2ac1e0475f1ff84d9438151cc9bd42967b52}
\end{flushleft} The 27 profile strings are the natural-language descriptions listed in Appendix~\ref{app:profiles}.

\subsection*{Low-burden scenario}

The low-burden condition used the following scenario text:

\begin{quote}
The person can schedule the vaccination appointment online. The appointment is available at a nearby pharmacy with extended evening and weekend hours. No documentation is required beyond a standard identification card. The appointment can be confirmed immediately, and the estimated total time required, including check-in and vaccination, is about 15 minutes.
\end{quote}

\subsection*{High-burden scenario}

The high-burden condition used the following scenario text:

\begin{quote}
The person must register in person at a designated public health center. The person must bring income verification and proof of residency. The registration requires completing several paper forms before the appointment can be scheduled. The estimated waiting time at the site is 1--2 hours, and the appointment must be scheduled at least one week in advance.
\end{quote}

\subsection*{Pairing of low- and high-burden conditions}

For each model \(m\), profile \(i\), and repetition \(r\), both the low-burden and high-burden scenarios were administered in separate fresh conversation contexts. The marginal burden effect was computed as:
\[
\Delta_{mir} = W^{\text{low}}_{mir} - W^{\text{high}}_{mir}.
\]
Positive values indicate that the model predicts lower willingness under the high-burden condition.

\subsection*{Burden-level comprehension diagnostic}

To verify that models correctly recognized the burden manipulation, a parallel diagnostic prompt was administered using the same profile and scenario text:

\begin{quote}
Individual profile:

\texttt{\{PROFILE\_TEXT\}}

Vaccination access scenario:

\texttt{\{SCENARIO\_TEXT\}}

Question:

Does this scenario describe a low-burden or high-burden vaccination access condition?

Return only one word: \texttt{low} or \texttt{high}.
\end{quote}

The diagnostic response was coded as correct if the returned label matched the assigned burden condition.

\section{Narrative Coding Codebook}
\label{app:narrative-codebook}

This appendix reports the algorithmic coding protocol used to classify high-burden rationales into five interpretive frames. The unit of coding is the rationale (one to three sentences) associated with each model--profile--repetition observation under the high-burden condition. Each frame is coded as a binary indicator using regular-expression pattern matching. Frames are not mutually exclusive.

\subsection{Algorithmic Coding Procedure}

Rationales are coded using a rule-based algorithm that applies predefined keyword patterns (regular expressions) to each rationale text. The algorithm operates as follows:

\begin{enumerate}
\item \textbf{Text preprocessing:} Rationale text is converted to lowercase.
\item \textbf{Pattern matching:} For each of the five frames, the algorithm searches for predefined regular expression patterns (word boundaries applied where appropriate).
\item \textbf{Binary assignment:} A frame is coded as present (1) if at least one pattern matches; absent (0) otherwise.
\item \textbf{Multiple frames:} Because frames are not mutually exclusive, a single rationale may be coded as positive for multiple frames.
\end{enumerate}

\subsection{Pattern Dictionary}

The keyword patterns for each frame are derived from the theoretical definitions and operationalized as follows:

{\raggedright
\begin{description}
\item[Autonomy infringement:] \texttt{individual choice}, \texttt{personal choice}, \texttt{personal freedom}, \texttt{autonomy}, \texttt{privacy}, \texttt{consent}, \texttt{personal discretion}, \texttt{right to choose}, \texttt{individual liberty}, \texttt{procedural dignity}, and phrases indicating limits on personal autonomy.
\item[Procedural legitimacy:] \texttt{justified}, \texttt{safeguard}, \texttt{reasonable verification}, \texttt{necessary procedure}, \texttt{eligibility}, \texttt{fraud prevention}, \texttt{legitimate public need}, \texttt{due process}, \texttt{accountability}, \texttt{oversight}, and phrases indicating procedural justification.
\item[Collective responsibility:] \texttt{civic duty}, \texttt{public health}, \texttt{protect others}, \texttt{communal benefit}, \texttt{vulnerable group}, \texttt{collective welfare}, \texttt{common good}, \texttt{herd immunity}, and phrases emphasizing community benefit.
\item[Access barriers:] \texttt{cannot afford}, \texttt{miss work}, \texttt{lose income}, \texttt{childcare}, \texttt{no regular doctor}, \texttt{lack of access}, \texttt{low income}, \texttt{unable to comply}, \texttt{disadvantage}, \texttt{single parent}, \texttt{hourly worker}, and phrases linking burdens to resource constraints.
\item[Coercive backlash:] \texttt{resentment}, \texttt{distrust}, \texttt{resistance}, \texttt{backlash}, \texttt{reduced legitimacy}, \texttt{mandate}, \texttt{penalty}, \texttt{enforcement}, \texttt{backfire}, \texttt{coercion}, and phrases indicating reactance to mandates.
\end{description}
}

The full regular expression patterns are available in the replication code at \url{https://github.com/yichao2022/llm-policy-simulation}, which is released under the MIT License (Copyright 2026 Yichao Jin); reuse and modification are permitted with attribution, and model outputs remain subject to the providers' terms of service.

\subsection{Validation Protocol}

To validate the algorithmic coding, a random sample of 200 rationales (approximately 10\% of the high-burden sample, $N \approx 2,019$) was manually inspected by the author. The validation protocol was:

\begin{enumerate}
\item \textbf{Blinded inspection:} The validator reviewed rationale text and algorithmic assignments without access to model identity or predicted outcomes.
\item \textbf{Definition check:} Each positive algorithmic assignment was checked against the frame definitions in Table~\ref{tab:narrative-codebook}.
\item \textbf{Focus areas:} Special attention was given to access-barrier coding, where generic mentions of paperwork or waiting time are not sufficient for a positive code under the algorithmic rules.
\item \textbf{Discrepancy review:} Any rationales where pattern matches appeared to violate frame definitions were flagged for review.
\end{enumerate}

The manual inspection found that pattern assignments reflected the codebook definitions rather than superficial keyword matches, and no systematic misclassification pattern was identified that would require pattern refinement. This review is a single-coder quality check on a 10\% sample; it establishes that the coding is transparent and rule-based, not that the frames are measured with validated reliability. The coded indicators are therefore used as descriptive inputs to the frame--outcome association analysis, and their coefficients should be read accordingly.

\subsection{Frame Prevalence and Exclusions}

\begin{table}[htbp]
\centering
\caption{Narrative coding codebook}
\label{tab:narrative-codebook}
\begin{tabular*}{\textwidth}{@{\extracolsep{\fill}}L{3.2cm}L{5.9cm}L{5.9cm}@{}}
\toprule
Frame & Code as 1 when rationale... & Do not code as 1 when rationale... \\
\midrule
Autonomy infringement &
Frames procedures as threats to individual choice, autonomy, privacy, personal freedom, consent, or procedural dignity. &
Mentions inconvenience, waiting, forms, or documentation without linking them to autonomy, privacy, dignity, or choice. \\

Procedural legitimacy &
Frames administrative requirements as justified, reasonable, legitimate, or necessary safeguards for verification, program integrity, safety, or orderly administration. &
Mentions documentation or registration without justifying it, or treats procedures only as obstacles. \\

Collective responsibility &
Emphasizes civic duty, public health responsibility, protection of others, community benefit, or willingness to tolerate inconvenience for collective welfare. &
Mentions only individual health benefits or convenience without reference to public health, community, or others. \\

Access barriers &
Explicitly links paperwork, waiting time, documentation, travel, scheduling, or in-person registration to reduced access, inability to comply, unequal burden, low income, work constraints, childcare, lack of regular provider access, or other vulnerability. &
Merely restates that the scenario involves forms, documentation, waiting time, travel, or delay without linking these requirements to access, resources, or vulnerability. \\

Coercive backlash &
Mentions resentment, distrust, resistance, backlash, reactance, reduced legitimacy, or negative reaction to mandates, penalties, enforcement, or perceived coercion. &
Describes practical inconvenience or access difficulty without resistance, distrust, resentment, or coercion. \\
\bottomrule
\end{tabular*}
\end{table}

The access-barrier definition was intentionally restricted to avoid coding simple restatements of the high-burden prompt as a substantive frame. For example, ``the process requires forms and a long wait'' is not coded as access barriers unless the rationale links those requirements to reduced access or profile-specific constraints such as work flexibility, childcare, income, transportation, or lack of regular care.

Frames with little or no empirical variation are reported descriptively but excluded from the main frame-adjusted model. In the present analysis, autonomy infringement and procedural legitimacy each appeared in one of the 2,019 high-burden rationales (0.05\% each), so neither frame has usable variation. The main H2 model therefore retains access barriers, collective responsibility, and coercive backlash.

\section{H2 Leave-One-Model-Out Robustness}
\label{app:h2-robustness}

To assess whether the H2 frame--outcome associations are driven by any single model, we re-estimated the fixed-effects frame model fifteen times, each time excluding one model from the expanded sample. The specification is the one used in the main text: model and profile fixed effects with the three retained frame indicators and no PVOC term. Table~\ref{tab:h2_leave_one_out} summarizes the minimum and maximum coefficient estimates across the leave-one-model-out specifications. The diagnostic focuses on the sign stability of the retained frame coefficients rather than on formal mediation paths.

\begin{table}[htbp]
\centering
\caption{Leave-one-model-out robustness for H2 frame coefficients}
\label{tab:h2_leave_one_out}
\resizebox{\textwidth}{!}{%
\begin{tabular}{lrrc}
\toprule
Variable & Minimum coefficient & Maximum coefficient & Sign stable \\
\midrule
Access barriers & +0.91 & +1.29 & Yes \\
Collective responsibility & -0.42 & +0.49 & No \\
Coercive backlash & +1.12 & +2.12 & Yes (never significant) \\
\bottomrule
\end{tabular}
}
\begin{flushleft}
\footnotesize Notes: Each row summarizes coefficients from fifteen leave-one-model-out specifications of the main H2 model (model and profile fixed effects, three frame indicators, no PVOC). The access-barrier coefficient is positive and significant in all fifteen specifications; the collective-responsibility and coercive-backlash coefficients, although positive for the latter, are not statistically distinguishable from zero in any specification. For comparison, a specification without fixed effects that instead includes the PVOC covariate yields access barriers 2.36 (\(p=0.019\)), collective responsibility 0.69 (\(p=0.515\)), and coercive backlash 7.27 (\(p<0.001\)), with a PVOC coefficient of \(-2.00\) (\(p<0.001\)); the gap is concentrated in the coercive-backlash contrast and indicates that this association is between-model rather than within-model.
\end{flushleft}
\end{table}

\section{Frame-Manipulation Prompts and H3 Robustness}
\label{app:h3-prompts-robustness}

This appendix reports the prompt templates and model-specific robustness checks for the within-model frame-manipulation experiment described in the main-text section on within-model frame manipulation and analyzed in the main-text H3 results. The purpose of the experiment is to test whether governance-value frames shift predicted administrative-burden effects while holding the base model, synthetic profiles, burden scenarios, output format, and decoding settings constant.

\subsection{Frame-Manipulation Prompt Templates}

The frame-manipulation experiment used four system-prompt frames. Each frame is built from the same two fixed sentences: an opening sentence and a closing instruction block, with the frame clause inserted between them. The neutral frame contains no frame clause, so the neutral system prompt is byte-identical to the baseline prompt reported in Appendix~\ref{app:burden-prompts}:
\begin{flushleft}
{\footnotesize\ttfamily SHA256 380d779e199673cdeaf808731aad01414b60e71edba7c23a57602241ff6dcffe}
\end{flushleft} the frame contrast is therefore measured against the true baseline condition. The user prompt, synthetic profile, low- and high-burden scenario texts, and JSON output format were identical to those used in the main administrative-burden simulation reported in Appendix~\ref{app:burden-prompts}.

\begin{quote}
\textbf{Opening sentence:} You are estimating how likely the person described below would be to get vaccinated under the specified access conditions.

\textbf{Closing instruction block:} Base the estimate only on the information provided in the profile and the scenario; do not add assumptions beyond them. Respond ONLY with a single JSON object and nothing else. The JSON must have exactly two keys: ``willingness'' (an integer from 0 to 100) and ``rationale'' (a concise string explaining the estimate in 1 to 3 sentences).
\end{quote}

The four frame conditions differ only in the inserted clause:

\begin{itemize}
    \item \textbf{Neutral frame:} no clause inserted.\\
    {\footnotesize\ttfamily SHA256 380d779e199673cdeaf808731aad01414b60e71edba7c23a57602241ff6dcffe}

    \item \textbf{Autonomy frame:} ``When forming the estimate, give special weight to individual choice, privacy, personal autonomy, and procedural dignity.''\\
    {\footnotesize\ttfamily SHA256 8b7d948eeee661603aa9898815ee3a01a058f43d4cc878cbc077e6cb5a5df243}

    \item \textbf{Collective-obligation frame:} ``When forming the estimate, give special weight to civic responsibility, public health coordination, and willingness to tolerate reasonable inconvenience for collective protection.''\\
    {\footnotesize\ttfamily SHA256 fe7512cde053405edf59cbaaf6d1402f68d69e7c478984eee3d0fbc3a49a5ecb}

    \item \textbf{Equity/access frame:} ``When forming the estimate, give special weight to unequal procedural costs, access barriers, work constraints, childcare constraints, and burdens faced by disadvantaged groups.''\\
    {\footnotesize\ttfamily SHA256 cba0156ababefd86611eb267ebf9ccddb1d60fe5a8cecf4e289871006042844a}
\end{itemize}

For each model \(m\), frame condition \(c\), profile \(i\), and repetition \(r\), the model evaluated both the low-burden and high-burden vaccination-access scenarios in fresh conversation contexts. The frame-specific burden effect was computed as:
\[
\Delta_{mcir} = W^{low}_{mcir} - W^{high}_{mcir}.
\]
Positive values indicate a larger predicted reduction in vaccination willingness under the high-burden condition.

\subsection{Mean Burden Effects by Model and Frame}

Table~\ref{tab:h3_model_frame_means} reports mean burden effects by model and frame condition for the four-model H3 subset. The descriptive means show that the collective-obligation frame lowers simulated burden effects in all four models, and that the autonomy frame increases simulated burden effects in three of the four models. The remaining contrasts are model-specific rather than uniform: for Qwen3.7 Plus the autonomy frame lowers the mean burden effect (53.71 vs.\ 56.59 under the neutral frame), and the equity/access frame has little effect for that model (55.64) and for Llama 3.1 70B Instruct (36.42 vs.\ 27.11 under neutral, a positive but smaller shift than the autonomy frame). The absolute level of the burden effect also differs substantially across models under the neutral frame (25.59 for GPT-4.1 to 56.59 for Qwen3.7 Plus), so frame effects are interpreted relative to each model's own neutral baseline.

Mistral Large is a particularly informative robustness case. It has the lowest PVOC score in the expanded fifteen-model sample, and all three frame manipulations shift its simulated burden effects in the expected directions. This suggests that frame sensitivity is not limited to models with high baseline PVOC or high baseline burden sensitivity.

\begin{table}[htbp]
\centering
\caption{Mean burden effects by model and frame}
\label{tab:h3_model_frame_means}
\small
\begin{tabular}{L{3.2cm}L{3.1cm}rrr}
\toprule
Model & Frame & Mean \(\Delta\) & SD & \(N\) \\
\midrule
GPT-4.1 & Neutral & 25.59 & 8.42 & 135 \\
GPT-4.1 & Autonomy & 34.15 & 6.84 & 135 \\
GPT-4.1 & Collective obligation & 18.62 & 6.22 & 135 \\
GPT-4.1 & Equity/access & 34.00 & 9.46 & 135 \\
\midrule
Llama 3.1 70B Instruct & Neutral & 27.11 & 9.69 & 135 \\
Llama 3.1 70B Instruct & Autonomy & 37.19 & 7.98 & 135 \\
Llama 3.1 70B Instruct & Collective obligation & 25.33 & 10.42 & 135 \\
Llama 3.1 70B Instruct & Equity/access & 36.42 & 12.90 & 134 \\
\midrule
Qwen3.7 Plus & Neutral & 56.59 & 7.64 & 135 \\
Qwen3.7 Plus & Autonomy & 53.71 & 6.54 & 135 \\
Qwen3.7 Plus & Collective obligation & 41.93 & 11.34 & 135 \\
Qwen3.7 Plus & Equity/access & 55.64 & 6.91 & 135 \\
\midrule
Mistral Large & Neutral & 37.04 & 6.44 & 135 \\
Mistral Large & Autonomy & 41.11 & 5.38 & 135 \\
Mistral Large & Collective obligation & 16.93 & 7.53 & 135 \\
Mistral Large & Equity/access & 39.63 & 6.95 & 135 \\
\bottomrule
\end{tabular}
\begin{flushleft}
\footnotesize Notes: Each row reports the mean profile--repetition-level burden effect, \(\Delta = W^{low}-W^{high}\), for one model under one assigned frame. $N=135$ for each model--frame cell, corresponding to 27 profiles and five repetitions, except where a response failed the structured-output checks: Llama 3.1 70B Instruct equity/access ($N=134$). Positive values indicate larger predicted reductions in willingness under the high-burden condition.
\end{flushleft}
\end{table}

\clearpage

\subsection{Small-Sample Inference for the Model-Specific Estimates}

Each model-specific regression rests on 27 profile clusters, so we checked the model-level conclusions against two small-sample procedures: a cluster covariance with the standard finite-sample correction ($G/(G-1)\cdot(N-1)/(N-K)$) and a wild cluster bootstrap-$t$ with Rademacher weights, the null imposed, and 1{,}999 replications. Table~\ref{tab:h3_small_sample} reports the resulting $p$-values next to the asymptotic cluster-robust ones. The three inferences agree: the marginal equity/access contrast for Mistral Large remains significant ($2.59$, $p=0.023$ asymptotic, $0.030$ corrected, $0.027$ bootstrap), the negative autonomy contrast for Qwen3.7 Plus remains significant ($-2.87$, $p=0.001$--$0.002$), and the collective-obligation contrast for Llama 3.1 70B Instruct and the equity/access contrast for Qwen3.7 Plus stay statistically indistinguishable from zero under every procedure.

\begin{table}[htbp]
\centering
\caption{Small-sample inference for model-specific H3 contrasts}
\label{tab:h3_small_sample}
\small
\resizebox{\textwidth}{!}{%
\begin{tabular}{llrrrr}
\toprule
Model & Contrast & Coefficient & $p$ (asymptotic) & $p$ (corrected) & $p$ (wild bootstrap) \\
\midrule
GPT-4.1 & Autonomy & $+8.56$ & $<0.001$ & $<0.001$ & $<0.001$ \\
GPT-4.1 & Collective obligation & $-6.97$ & $<0.001$ & $<0.001$ & $<0.001$ \\
GPT-4.1 & Equity/access & $+8.41$ & $<0.001$ & $<0.001$ & $<0.001$ \\
Llama 3.1 70B Instruct & Autonomy & $+10.07$ & $<0.001$ & $<0.001$ & $<0.001$ \\
Llama 3.1 70B Instruct & Collective obligation & $-1.78$ & 0.137 & 0.155 & 0.125 \\
Llama 3.1 70B Instruct & Equity/access & $+9.36$ & $<0.001$ & $<0.001$ & $<0.001$ \\
Qwen3.7 Plus & Autonomy & $-2.87$ & $<0.001$ & 0.001 & 0.002 \\
Qwen3.7 Plus & Collective obligation & $-14.65$ & $<0.001$ & $<0.001$ & $<0.001$ \\
Qwen3.7 Plus & Equity/access & $-0.94$ & 0.250 & 0.271 & 0.265 \\
Mistral Large & Autonomy & $+4.07$ & $<0.001$ & $<0.001$ & $<0.001$ \\
Mistral Large & Collective obligation & $-20.11$ & $<0.001$ & $<0.001$ & $<0.001$ \\
Mistral Large & Equity/access & $+2.59$ & 0.023 & 0.030 & 0.027 \\
\bottomrule
\end{tabular}
}
\begin{flushleft}
\footnotesize Notes: Each model-specific regression includes profile fixed effects and frame indicators, with standard errors clustered by profile (27 clusters per model). Asymptotic $p$-values use the cluster-robust covariance without small-sample correction; corrected $p$-values apply $G/(G-1)\cdot(N-1)/(N-K)$; wild bootstrap $p$-values come from 1{,}999 Rademacher-weighted bootstrap replications with the null imposed on the tested contrast. Bootstrap $p$-values are reported to three decimals and are bounded below by $1/2{,}000$.
\end{flushleft}
\end{table}

\subsection{Model-Specific Frame-Manipulation Effects}

Table~\ref{tab:h3_model_specific_effects} reports model-specific frame regressions for the four-model H3 subset. Each model-specific regression estimates \(\Delta\) as a function of frame indicators and profile fixed effects, using the neutral frame as the reference condition. The pooled column corresponds to the main H3 specification reported in the main text, which pools the same four models and includes model and profile fixed effects.

\begin{table}[htbp]
\centering
\caption{Model-specific frame-manipulation effects}
\label{tab:h3_model_specific_effects}
\small
\begin{tabular}{L{3.0cm}rrrrr}
\toprule
Frame contrast & GPT-4.1 & Llama 3.1 70B & Qwen3.7 Plus & Mistral Large & Pooled \\
\midrule
Autonomy vs. neutral 
& 8.56*** 
& 10.07*** 
& -2.87** 
& 4.07*** 
& 4.96*** \\
Collective obligation vs. neutral 
& -6.97*** 
& -1.78 
& -14.65*** 
& -20.11*** 
& -10.88*** \\
Equity/access vs. neutral 
& 8.41*** 
& 9.36*** 
& -0.94 
& 2.59* 
& 4.84*** \\
\bottomrule
\end{tabular}
\begin{flushleft}
\footnotesize Notes: Each model-specific column reports a regression of \(\Delta\) on frame indicators and profile fixed effects for the indicated model. The pooled column reports the main H3 specification for the four-model H3 subset with model and profile fixed effects. Standard errors are clustered at the model--profile level in all columns. The neutral frame is the omitted reference category. Positive coefficients indicate larger simulated burden effects relative to neutral; negative coefficients indicate smaller simulated burden effects. Standard errors and 95\% confidence intervals are reported in Table~\ref{tab:h3_model_specific_full}. \(^{*}p<0.05\), \(^{***}p<0.001\).
\end{flushleft}
\end{table}

Table~\ref{tab:h3_model_specific_full} reports the corresponding standard errors and confidence intervals. The model-specific results show both robustness and heterogeneity within the H3 subset. The collective-obligation frame reduces simulated burden effects in all four models (significantly in GPT-4.1, Qwen3.7 Plus, and Mistral Large; not statistically distinguishable from zero in Llama 3.1 70B Instruct). The autonomy frame increases simulated burden effects in GPT-4.1, Llama 3.1 70B Instruct, and Mistral Large, but \emph{reverses sign} for Qwen3.7 Plus ($-2.87$, 95\% CI $[-4.53, -1.22]$), where activating autonomy concerns reduces the simulated burden effect rather than amplifying it. The equity/access frame is positive and significant for GPT-4.1, Llama 3.1 70B Instruct, and Mistral Large, and statistically indistinguishable from zero for Qwen3.7 Plus. Mistral Large is especially informative because it has the lowest PVOC score in the expanded fifteen-model sample, yet all three frame manipulations shift its simulated burden effects in the expected directions. These results support the pooled H3 finding for the collective-obligation contrast while showing that model endpoints differ in both the magnitude and the sign of responsiveness to specific governance-value frames; the pooled autonomy and equity estimates are averages over a genuinely mixed set of model-level responses, and the Qwen3.7 Plus reversal is the clearest caution against reading a pooled frame effect as a property of LLMs in general.

\begin{table}[htbp]
\centering
\caption{Model-specific frame-manipulation effects with standard errors and confidence intervals}
\label{tab:h3_model_specific_full}
\small
\begin{tabular}{L{2.8cm}L{3.7cm}rrr}
\toprule
Model & Frame contrast & Coefficient & SE & 95\% CI \\
\midrule
GPT-4.1 & Autonomy vs. neutral & 8.56*** & 0.76 & [6.99, 10.12] \\
GPT-4.1 & Collective obligation vs. neutral & -6.97*** & 0.81 & [-8.63, -5.31] \\
GPT-4.1 & Equity/access vs. neutral & 8.41*** & 0.77 & [6.83, 9.99] \\
\midrule
Llama 3.1 70B Instruct & Autonomy vs. neutral & 10.07*** & 1.46 & [7.07, 13.08] \\
Llama 3.1 70B Instruct & Collective obligation vs. neutral & -1.78 & 1.21 & [-4.27, 0.72] \\
Llama 3.1 70B Instruct & Equity/access vs. neutral & 9.36*** & 1.75 & [5.76, 12.97] \\
\midrule
Qwen3.7 Plus & Autonomy vs. neutral & -2.87** & 0.81 & [-4.53, -1.22] \\
Qwen3.7 Plus & Collective obligation vs. neutral & -14.65*** & 1.06 & [-16.84, -12.47] \\
Qwen3.7 Plus & Equity/access vs. neutral & -0.94 & 0.84 & [-2.66, 0.78] \\
\midrule
Mistral Large & Autonomy vs. neutral & 4.07*** & 0.98 & [2.05, 6.09] \\
Mistral Large & Collective obligation vs. neutral & -20.11*** & 0.89 & [-21.93, -18.29] \\
Mistral Large & Equity/access vs. neutral & 2.59* & 1.13 & [0.28, 4.91] \\
\midrule
Pooled & Autonomy vs. neutral & 4.96*** & 0.70 & [3.57, 6.35] \\
Pooled & Collective obligation vs. neutral & -10.88*** & 0.84 & [-12.54, -9.21] \\
Pooled & Equity/access vs. neutral & 4.84*** & 0.70 & [3.45, 6.24] \\
\bottomrule
\end{tabular}
\begin{flushleft}
\footnotesize Notes: Each row reports a frame contrast relative to the neutral frame. Model-specific regressions include profile fixed effects. The pooled regression includes model and profile fixed effects. Standard errors are clustered at the model--profile level for the pooled column (108 clusters) and at the profile level within each model for the model-specific columns (27 clusters each), because profiles are the controlled test cases and queries within a profile are repeated measures. The pooled specification uses $N=2{,}159$ profile--repetition observations (540 for GPT-4.1, 539 for Llama 3.1 70B Instruct, 540 for Qwen3.7 Plus, and 540 for Mistral Large, after excluding responses that failed the structured-output checks). The dependent variable is \(\Delta = W^{low}-W^{high}\). Positive coefficients indicate larger simulated burden effects relative to neutral. Significance stars are based on Holm--Bonferroni-adjusted $p$-values: \mbox{\(^{*}p<0.05\)}, \mbox{\(^{***}p<0.001\)}.
\end{flushleft}
\end{table}

\section{Additional Robustness Checks}
\label{app:additional-robustness}

This appendix reports additional robustness checks for the H1 orientation--effect analysis. The main H1 specification uses the composite PVOC index as the model-level predictor of average burden effects. Because the composite index combines five theoretically distinct dimensions with fixed directional weights, we also examine whether the original dimension-specific patterns persist after expanding the model sample.

\subsection{Caveats Qualified in the Main Text}

Three caveats qualify the null H1 result. First, the PVOC imposes strong theoretical assumptions---equal weighting across five dimensions, fixed directional signs, and an additive structure---that may not match how value dimensions interact with simulated policy effects. Second, internal consistency is mixed: only collective obligation reaches conventional reliability ($\alpha = 0.869$), so a composite of components that do not cohere as a measurement scale cannot be expected to predict an external outcome. Third, at $N = 15$ model endpoints, model-level regressions have limited power, which is why the power and minimum-detectable-effect diagnostics below are reported alongside the coefficient estimates.

\subsection{Expanded-Sample H1 Robustness}

Table~\ref{tab:h1_expanded_robustness} summarizes the key H1 robustness results for the original nine-model sample and the expanded fifteen-model sample. The composite PVOC index does not predict model-level burden effects in either sample. In the original nine-model sample, the PVOC coefficient is negative and statistically indistinguishable from zero (\(\hat{\beta}=$-$4.24\), \(p=0.219\), \(R^2=0.206\)). After expanding to fifteen endpoints, the PVOC coefficient remains weak and statistically indistinguishable from a null effect at conventional power (\(\hat{\beta}=$-$2.14\), \(R^2=0.090\)).

The expanded sample also changes the interpretation of the dimension-specific collective-obligation result. In the original nine-model sample, collective obligation (CO) was positively associated with model-level burden effects but did not reach conventional significance (\(\hat{\beta}=10.81\), \(p=0.179\), \(R^2=0.242\)). After the six additional endpoints are added, the association shrinks and remains non-significant (\(\hat{\beta}=1.28\), \(p=0.725\), \(R^2=0.010\)). This indicates that the CO association is sensitive to model composition and should be interpreted as an exploratory signal at best rather than as a stable dimension-level predictor.

The expanded-sample measurement diagnostics also suggest caution in interpreting the composite PVOC formula. CO is the only dimension with acceptable exploratory internal consistency, but removing CO from the composite does not change the substantive conclusion: the leave-CO-out PVOC coefficient remains statistically indistinguishable from zero (Table~\ref{tab:h1_expanded_robustness}). This does not validate an alternative index; the composite association is weak under either weighting.

\begin{table}[htbp]
\centering
\caption{H1 robustness across original and expanded model samples}
\label{tab:h1_expanded_robustness}
\begin{tabular}{lrrrl}
\toprule
Predictor and sample & Coefficient & \(p\) & \(R^2\) & Interpretation \\
\midrule
PVOC composite, \(N=9\) & $-$4.24 & 0.219 & 0.206 & Null \\
PVOC composite, \(N=15\) & $-$2.14 & 0.278 & 0.090 & Null \\
Collective obligation, \(N=9\) & 10.81 & 0.179 & 0.242 & Not significant \\
Collective obligation, \(N=15\) & 1.28 & 0.725 & 0.010 & Not stable \\
Leave-CO-out PVOC, \(N=15\) & $-$2.22 & 0.333 & 0.072 & Null \\
\bottomrule
\end{tabular}
\begin{flushleft}
\footnotesize Notes: Each row reports a separate bivariate model-level regression of \(\bar{\Delta}_m\) on the listed predictor or PVOC variant. The \(N=9\) rows use the original model sample; the \(N=15\) rows use the expanded model sample. Higher \(\bar{\Delta}_m\) indicates a larger predicted reduction in willingness under high administrative burden. The leave-CO-out PVOC result is reported as a sensitivity diagnostic, not as a validated alternative index.
\end{flushleft}
\end{table}

\subsection{Leave-One-Model-Out Diagnostics}

Leave-one-model-out diagnostics further support the main H1 conclusion. The composite PVOC coefficient is not a stable predictor of model-level burden effects. In the expanded sample, the bivariate PVOC association remains weak and statistically indistinguishable from zero. The CO pattern is not significant in either sample: the coefficient is positive but small in the nine-model analysis and shrinks further after adding six additional endpoints.

These diagnostics reinforce the interpretation that H1 does not provide evidence for a stable aggregate orientation--effect relationship. The PVOC index is useful as a theoretically motivated diagnostic construct, but the expanded analysis does not support treating it as a validated scalar predictor of simulated burden-effect magnitude.

\subsection{Leave-One-Profile-Out Diagnostics}

We re-estimated each of the three main models twenty-seven times, excluding one synthetic profile at a time (27 fits per model). Coefficient ranges across the 27 fits are archived with the replication package. The H1 PVOC slope stays between $-2.26$ and $-2.07$ and is never statistically distinguishable from zero ($p$ between 0.26 and 0.30), so the H1 null is not produced by any single test case. For H2, the access-barrier coefficient remains positive and significant in all 27 fits (range $+0.89$ to $+1.26$, $p<0.05$ throughout); the collective-responsibility coefficient is indistinguishable from zero in every fit ($-0.27$ to $+0.35$), and the coercive-backlash coefficient, although positive in every fit ($+1.10$ to $+2.12$), is never statistically significant. For H3, all three frame contrasts retain both sign and significance in all 27 fits (autonomy $+4.64$ to $+5.18$; collective obligation $-11.01$ to $-10.70$; equity/access $+4.56$ to $+5.10$; $p<0.001$ throughout).

\subsection{Power and Minimum Detectable Effect Diagnostics}
\label{app:h1-power}

The original nine-model H1 analysis was underpowered for small or moderate model-level associations. In that sample, the approximate 80\% minimum detectable slope for the composite PVOC index was \(\beta \approx 10.2\) (5.7 in the fifteen-endpoint sample). The observed PVOC coefficient was far smaller than this threshold. Expanding the model sample to fifteen endpoints improves coverage but does not eliminate the basic limitation that H1 relies on model-level variation.

For this reason, the H1 estimates should be interpreted as descriptive diagnostics rather than high-powered population estimates. The expanded results are most informative as evidence against a large linear association between the composite PVOC index and simulated burden effects in this model sample. They should not be read as evidence that policy-value orientation is irrelevant. Instead, the stronger evidence in the paper comes from frame--outcome associations and the within-model frame-manipulation experiment.

\end{document}
